\documentclass[10pt,a4paper]{amsart}
\usepackage[margin=19mm]{geometry}
\usepackage{amsmath,amssymb,mathtools}
\usepackage[scr=rsfs,cal=euler]{mathalfa}
\usepackage{xfrac}
\usepackage[unicode,hidelinks,bookmarks=false]{hyperref}
\newcommand{\ie}{i.e.\ }
\newcommand{\cf}{cf.\ }
\newcommand{\resp}{respectively\ }
\newcommand{\Subsection}{\subsection}
\providecommand{\nobreakdash}{\nobreak\hbox{-}}
\setlength{\emergencystretch}{2em}
\multlinegap0pt
\usepackage{amssymb}
\usepackage{mathtools}
\usepackage{bm}
\usepackage[shortlabels]{enumitem}
\newtheorem{lemma}{Lemma}
\newcommand{\BB}{\bm B}
\newcommand{\CC}{\bm C}
\newcommand{\HH}{\bm H}
\newcommand{\Id}{\bm I}
\newcommand{\calD}{\mathcal D}
\newcommand{\dd}{\,\mathrm d}
\newcommand{\nn}{\bm n}
\newcommand{\tr}{\operatorname{tr}}
\title{Entropy-Compatible Barrier Schemes for Diffusive FENE Flows}
\author{Sai Peng}
\date{Source: arXiv:2606.05209v1 (CC BY 4.0)}
\begin{document}
\maketitle

\noindent\emph{Source and licence.} Entropy-Compatible Barrier Schemes for Diffusive FENE Flows. Version arXiv:2606.05209v1 (CC BY 4.0), distributed by arXiv under the Creative Commons Attribution 4.0 International licence, http://creativecommons.org/licenses/by/4.0/, as recorded in the arXiv licence metadata for this exact version and shown on the abstract page https://arxiv.org/abs/2606.05209v1. Full licence notice: \texttt{licenses/arxiv-2606.05209-CC-BY-4.0.txt}.\par\smallskip

\noindent\textit{Editorial scope.} Counted unit: the article's lemma lemma (source0.tex lines 265--283). The article's complete original proof of it is delivered with it (source0.tex lines 285--294, one \texttt{proof} environment of 10 lines). No mathematics is written, repaired or re-derived: the statement and the proof are the article's own lines, in the article's own notation, and no retained line is substituted or re-flowed.  Retained uncounted as the source-local context the proof needs: lines 197--220.  Nothing was omitted from any retained span, so the package carries no omission record.  No retained line contains a citation, so the wrapper carries no bibliography and no bibliography entry.  No retained line contains a figure environment, an image-inclusion command or drawing code, so no image is delivered and the package declares no asset dependency.  Editorial formatting only, in addition: an AMS \texttt{amsart} preamble that restates the article's own declarations and macros used by the retained lines, namely the environments \texttt{lemma} on separate counters, together with the standard packages those lines need.  The retained environments print in this document as Lemma 1 in order of appearance, whereas the article prints the same items with the numbering of its own full document.  The complete native source of the article as distributed in the arXiv e-print for this exact version is bundled unchanged as \texttt{sources/202100/source0.tex}.
The elastic entropy is
\[
  \Phi_b(\CC)
  =
  -\log\det\CC-(b-d)\log\left(\frac{b-\tr\CC}{b-d}\right).
  \label{eq:fene-entropy}
\]
It satisfies $\Phi_b(\Id)=0$ and
\[
  D\Phi_b(\CC)[\BB]
  =
  \left(f_b(\CC)\Id-\CC^{-1}\right):\BB
  =\HH_b(\CC):\BB .
\]
The Hessian is positive:
\[
  D^2\Phi_b(\CC)[\BB,\BB]
  =
  \tr(\CC^{-1}\BB\CC^{-1}\BB)
  +
  \frac{b-d}{(b-\tr\CC)^2}\bigl(\tr\BB\bigr)^2 .
  \label{eq:hessian}
\]
Thus $\Phi_b$ is convex on $\calD_b$ and blows up at both parts of the boundary: $\lambda_{\min}(\CC)\downarrow0$ and $\tr\CC\uparrow b$.

\begin{lemma}[polymer molecular-diffusion entropy dissipation]
Assume periodic boundary conditions or $\nabla\CC\,\nn=0$ on $\partial\Omega$.  For every smooth admissible conformation field,
\[
  \int_\Omega \Delta\CC:\HH_b(\CC)\,\dd x
  =
  -\int_\Omega \mathcal I_b(\CC)\,\dd x ,
\]
where
\[
  \mathcal I_b(\CC)
  =
  \sum_{j=1}^d
  \left[
  \tr\!\left(\CC^{-1}\partial_j\CC\,\CC^{-1}\partial_j\CC\right)
  +
  \frac{b-d}{(b-\tr\CC)^2}\bigl(\partial_j\tr\CC\bigr)^2
  \right]\ge0 .
\]
\end{lemma}

\begin{proof}
Integrating by parts gives
\[
  \int_\Omega \Delta\CC:\HH_b(\CC)\,\dd x
  =
  -\sum_{j=1}^d\int_\Omega
  \partial_j\CC:D\HH_b(\CC)[\partial_j\CC]\,\dd x .
\]
Since $D\HH_b=D^2\Phi_b$, the Hessian formula \eqref{eq:hessian} gives the stated expression.
\end{proof}

\end{document}
